\section{Introduction}
\label{sec:intro}

The lid-driven cavity flow is a classical test case for incompressible viscous-flow solvers and a canonical configuration for studying confined recirculating motion. For the square geometry, benchmark velocity profiles have been established by Ghia, Ghia, and Shin~\cite{Ghia1982}, and spectral computations have subsequently provided highly accurate reference fields at low and moderate Reynolds numbers~\cite{Botella1998,Albensoeder2005}. The review of Shankar and Deshpande~\cite{Shankar2000} summarizes the richness of the configuration, which combines a simple boundary geometry with corner eddies, non-uniqueness, and transition to unsteadiness. Early analytical and numerical studies clarified the internal recirculation and the corner-eddy structure~\cite{Batchelor1956,Burggraf1966,Pan1967}, while the local analysis of Moffatt~\cite{Moffatt1964} established that a sequence of viscous eddies forms near a sharp corner and that the local stresses grow as the corner is approached. These results make the lid-driven cavity a natural setting in which to ask how a controlled modification of the wall shape redistributes the global energy budget.

The lattice Boltzmann method (LBM) has become a standard kinetic approach for such flows~\cite{Benzi1992,Chen1998,Succi2001}, and it has been applied to the square cavity since the work of Hou \emph{et al.}~\cite{Hou1995}. Deep and rectangular cavities have also been studied with LBM~\cite{Patil2006}, showing that the method remains accurate when the recirculation is strongly anisotropic. For the present study, three ingredients of the kinetic formulation matter. First, a multiple-relaxation-time (MRT) collision operator is used, which separates the relaxation of hydrodynamic and non-hydrodynamic moments and improves stability at the Reynolds numbers considered here~\cite{dHumieres2002,Lallemand2000}. Second, the curved bottom wall is treated with interpolated bounce-back~\cite{Bouzidi2001}, while the moving lid uses a link-based velocity closure of the type reviewed by Latt \emph{et al.}~\cite{Latt2008}. Third, the momentum forcing and the velocity boundary conditions are introduced consistently with the discrete lattice effects analyzed by Guo \emph{et al.}~\cite{Guo2002} and with the pressure--velocity closures of Zou and He~\cite{Zou1997}. These choices are not incidental: the energetic quantities analyzed below are volume integrals of the stress and velocity fields, so their accuracy depends directly on how the near-wall populations are reconstructed.

Geometric variations involving corrugated or wavy bottom walls have received particular attention. Bisht, Haeri, and Patil~\cite{Bisht2017} investigated wall-driven enclosures with regular square corrugations using LBM and a fictitious-domain method, focusing on flow topology and on the influence of corrugation height and number. More recently, Fatima \emph{et al.}~\cite{Fatima2023} conducted D2Q9 lattice Boltzmann simulations in a lid-driven cavity with corrugated walls of various wavelengths, focusing on the resulting recirculation patterns, secondary corner vortices, and streamline topology. Mixed-convection studies with sinusoidal wavy bottoms~\cite{AlAmiri2007,Nasrin2011} have examined heat transfer in similar enclosures and have shown that the wall wavelength and amplitude modify the recirculation and the thermal field. Wavy cavities of the nanofluid and mixed-convection type have likewise been examined for heat-transfer augmentation~\cite{AbuNada2014}. In those works, however, the primary focus has been on streamline topology, recirculation zones, and vortex formation. Global mechanical energetics---the integrated power input, the volume-integrated viscous dissipation, matched pairwise contrasts, and their systematic comparison with an area-based extensive reference---were not evaluated. The literature therefore demonstrates that corrugated walls alter recirculating patterns, but it has not yet provided a systematic energetic accounting of how successive, controlled changes in the bottom-wall geometry redistribute the power input and viscous dissipation relative to what would be expected from a simple change in fluid area alone.

A distinctive feature of the present configuration is that the bottom geometry is not merely corrugated: it is a \emph{superelliptic} family whose exponent $n$ continuously interpolates between a semicircular bottom ($n=2$, exact for the width-to-depth ratio used here) and a near-square cavity ($n\to\infty$), on which a single controlled undulation of amplitude $\alpha$ is superimposed. This makes the family a clean one-parameter probe of the corner geometry at fixed cavity width and lid velocity, in the spirit of the corner-singularity analyses of Moffatt~\cite{Moffatt1964} and of the near-lid energy concentration documented for the cubical cavity by Leriche and Gavrilakis~\cite{Leriche2000}. The key distinction from prior corrugated-cavity investigations is three-fold:
first, rather than documenting streamline topology and vortex centers qualitatively, the present study quantifies matched pairwise contrasts of the global energetic integrals ($P^*$, $D^*$, $K^*$);
second, it introduces an exact $\mathrm{d}A/\mathrm{d}B$ spatial partition that separates the dissipation contribution of the newly opened shoulder region from the redistribution within the common cavity core---this decomposition reveals the physical mechanism underlying the sub-extensive response;
and third, an independent conforming Taylor--Hood finite-element solver provides a cross-validation of the primary energetic response.

The central original contribution of this work is the exact $\mathrm{d}A/\mathrm{d}B$ spatial partition of the dissipation contrast. While spatial decompositions are common in mathematical analysis, their application to the energetic accounting of cavity flows---separating the added-region dissipation from the common-region redistribution---provides direct physical insight into why the dissipation response is sub-extensive relative to the area change. This partition demonstrates that the added shoulder region, despite occupying $7.4\%$ of the fluid area at $Re=400$, contributes only $0.2\%$ of the total dissipation, while the common-region redistribution is itself a near-cancellation of opposing lid-band and core-band contributions.

The configuration also differs from the curved-wall cavity studies that are closest to it. Cavities with a purely semicircular bottom have been analyzed for their stability and transition: Pan, Li, and Chiu~\cite{Pan2025} computed the critical Reynolds number for the onset of unsteadiness in a three-dimensional semicircular-bottom cavity and characterized the oscillatory mode; the $n=2$ member of the present family is exactly a two-dimensional semicircular bottom of radius equal to half the cavity width, so it corresponds to the two-dimensional cross-section of that geometry. The present study works strictly below such thresholds and focuses on the \emph{steady} mechanical-energy budget rather than on bifurcation. The family is designed so that increasing $n$ removes curvature continuously and introduces progressively sharper shoulder corners---precisely the feature that drives the energetic response analyzed here. On the numerical side, high-Reynolds-number two-dimensional cavity flow has been computed with the same MRT-lattice-Boltzmann family of schemes used here by Chai, Shi, and Zheng~\cite{Chai2006}, which supports the stability of the present solver at the Reynolds numbers considered, but those works consider the classical square cavity and do not address geometric contrasts or energetic accounting. The contribution of the present work is to combine a continuous superelliptic interpolation of the bottom wall with (i) matched geometric contrasts of the integrated power and dissipation, (ii) a comparison against an area-only extensive reference, and (iii) the exact $\mathrm{d}A/\mathrm{d}B$ partition of the dissipation difference into its common-region and added-region contributions. This combination provides a complete energetic accounting that was not previously available for corrugated or curved-wall cavity flows.

The images of the flow field alone do not fully determine the energetic cost of the motion. Two configurations can exhibit markedly different streamline topologies without a proportionate difference in the integrated mechanical power, and conversely configurations with similar-looking flow patterns can present measurable differences in the volume-integrated viscous dissipation. Number, size, and position of recirculation zones are therefore useful qualitative diagnostics, but they do not by themselves quantify the global energy budget or its sensitivity to the geometry. This limitation motivates an energetic analysis based on integrated quantities and on matched, pairwise comparisons between configurations that differ in a controlled way in a single geometric parameter. It also motivates a careful separation between the convergence of absolute integrals and the stability of their differences; the standard framework for discretization-uncertainty estimation~\cite{Celik2008} motivates reporting this distinction explicitly rather than assuming grid independence.

The present work investigates how controlled changes in the bottom-wall geometry affect integrated energetic quantities in a lid-driven cavity, using matched geometric contrasts and a separate assessment of numerical consistency. Rather than introducing a new universal scaling law for corrugated cavities, the contribution is a quantitative, resolution-aware assessment built on three elements: (i) matched contrasts between successive superellipse exponents $n=2\rightarrow4\rightarrow6\rightarrow8$; (ii) comparison with an area-only extensive reference; (iii) independent evaluation of power-based and dissipation-based contrasts across $N=128,256,384$. This approach allows separation between the robustness of \emph{differences} between paired cases and the convergence of \emph{absolute} energetic integrals, which are distinct numerical questions.

The objectives are: (1) to quantify the influence of $(n,\alpha)$ on $P^*$, $D^*$, and $K^*$ over $Re\in\{100,400,1000\}$; (2) to measure matched contrasts $I_Q(2,4)$, $I_Q(4,6)$, and $I_Q(6,8)$ and test them against an area-only extensive reference; (3) to localize the origin of the dissipation contrast through a layer decomposition of the local strain-rate field and to connect it to the corner-flow mechanism; (4) to examine the Reynolds-number dependence and the effect of the undulation amplitude $\alpha$ on the hierarchy of these contrasts; (5) to assess the robustness of the contrasts under changes of wall treatment, of the effective Reynolds number, and of grid resolution, with explicit separation between estimator consistency ($I_P/I_D$) and quantitative convergence of $I(2,4)$ itself. The manuscript proceeds as follows: Section~\ref{sec:formulation} defines the geometry, the governing equations, and the energetic measures; Section~\ref{sec:numerical} describes the MRT-lattice Boltzmann method, wall treatments, and the numerical parameter matrix; Section~\ref{sec:results} presents the results, including flow visualization, the spatial mechanism, the absolute energetic response, the reference models, the contrast hierarchy, the Reynolds dependence, the robustness tests, and consistency checks; and Section~\ref{sec:conclusions} summarizes the main findings and states the scope and limitations of the conclusions.
